\documentclass[11pt]{article}

\usepackage[T1]{fontenc}
\usepackage{amsmath,amssymb,amsthm}
\usepackage[margin=1in]{geometry}
\usepackage{booktabs}
\usepackage{microtype}

\newcommand{\bits}{\{0,1\}}
\newcommand{\LC}{\operatorname{LC}}
\newcommand{\val}{\operatorname{val}}
\newcommand{\supp}{\operatorname{supp}}

\newtheorem{theorem}{Theorem}

\title{Explicit Hard Distributions, Boolean Tilts, and the Circuit Frontier}
\author{}
\date{}

\begin{document}

\maketitle

\begin{abstract}
This note unpacks the explicit full-support rational distribution used in the
Lean-checked exponential lower bound for cubic localization.  The distribution
itself is elementary: order the points of a Boolean cube by their binary value,
give them rapidly growing powers-of-powers weights, and normalize.  The point
of the unusual weights is that products and then sums of products carry two
independent layers of exact binary encoding.  The same collision can instead
be separated by a two-level Boolean tilt, yielding an existential exponential
NAND lower bound.  A structured block-parity subfamily goes further: Lean now
constructs canonical finite-search witnesses giving a bounded-precision,
quadratic-in-dimension cubic-localization lower bound.  The witnesses are
astronomically expensive to find, so the usual efficient-explicitness barrier
remains.
\end{abstract}

\section{The definition}

Fix a natural number $m$ and put
\[
  n_m := 4m+24.
\]
Thus $D_m$ is a distribution on $n_m$ visible bits.  For
$x=(x_0,\ldots,x_{n_m-1})\in\bits^{n_m}$, use the little-endian binary value
\[
  \val(x):=\sum_{i=0}^{n_m-1}x_i2^i
  \quad\in\quad\{0,\ldots,2^{n_m}-1\}.
\]
Let
\[
  f:=(1,\ldots,1),
  \qquad \val(f)=2^{n_m}-1,
\]
be the distinguished \emph{filler state}.  Define the unnormalized integer
weight
\[
  W_m(x):=
  \begin{cases}
    1, & x=f,\\[2mm]
    2^{\,2^{\val(x)}}, & x\ne f.
  \end{cases}
\]
Finally, let
\[
  Z_m:=\sum_{y\in\bits^{n_m}}W_m(y),
  \qquad
  D_m(x):=\frac{W_m(x)}{Z_m}.
\]
Equivalently,
\[
  Z_m=1+\sum_{v=0}^{2^{n_m}-2}2^{2^v}.
\]
This is the whole definition.  Every $W_m(x)$ is a positive integer and
$Z_m$ is a positive integer, so $D_m$ is an exact rational distribution with
\[
  \supp(D_m)=\bits^{n_m}.
\]
The index $m$ is only a scale parameter: it determines the dimension
$n_m=4m+24$.  It is not a hidden variable and it does not index a mixture.

\section{A tiny version of the table}

The theorem starts at $n=24$, so its actual tables are already enormous.  To
see the pattern, apply the same recipe on two bits.  Here
$\val(x_0,x_1)=x_0+2x_1$, and the table is
\[
\begin{array}{ccccc}
\toprule
x_0 & x_1 & \val(x) & W(x) & D(x)\\
\midrule
0&0&0&2^{2^0}=2&2/23\\
1&0&1&2^{2^1}=4&4/23\\
0&1&2&2^{2^2}=16&16/23\\
1&1&3&1&1/23\\
\bottomrule
\end{array}
\]
The non-filler weights begin
\[
  2,\ 4,\ 16,\ 256,\ 65536,\ldots,
\]
and satisfy $W(v+1)=W(v)^2$.  The all-ones state is deliberately reset to
weight one so that it can be used to pad products without changing their
unnormalized value.

The distribution is consequently extremely skewed: almost all of its mass is
on the last non-filler state.  Nevertheless every point has positive mass, so
the lower bound comes entirely from the exact relative weights, not from a
complicated support.

\section{Why these weights are useful}

Write
\[
  B:=2^n-1
\]
for the number of non-filler visible states in an $n$-bit version of the
construction.  Index those states by $v=0,\ldots,B-1$.  A candidate is a
binary vector
\[
  c=(c_0,\ldots,c_{B-1})\in\bits^B,
\]
or equivalently a subset of the non-filler states.  Its ordinary binary code
is
\[
  \operatorname{code}(c):=\sum_{v=0}^{B-1}c_v2^v.
\]

The formal construction associates to $c$ a homogeneous visible monomial of
degree $2B$.  It uses one copy of state $v$ when $c_v=1$ and fills all unused
positions with copies of $f$.  Since $W(f)=1$, its unnormalized value is
\[
  \prod_{v:c_v=1} W(v)
  =\prod_{v:c_v=1}2^{2^v}
  =2^{\sum_v c_v2^v}
  =2^{\operatorname{code}(c)}.
\]
After normalization, every such monomial has the common extra factor
$Z^{-2B}$.

There are now two exact uniqueness mechanisms.
\begin{enumerate}
  \item Different candidates $c$ have different values
  $\operatorname{code}(c)$, by uniqueness of binary expansion.
  \item If $A$ is a set of candidates, then
  \[
    \sum_{c\in A}2^{\operatorname{code}(c)}
  \]
  uniquely determines $A$, again by uniqueness of binary expansion.
\end{enumerate}
Thus distinct sets of candidate monomials always have different values on
$D_m$.  This nested coding is the reason for the nested exponent
$2^{2^{\val(x)}}$.

\section{Where cubic locality enters}

Suppose, toward a contradiction, that $D_m$ were the visible marginal of a
3-local law with $\ell$ hidden bits.  The joint system then has $n+\ell$
binary coordinates.  A degree-at-most-three feature is indexed by a scope of
at most three coordinates, and the Lean development proves the convenient
bound
\[
  \#\{\text{cubic scopes}\}\le (n+\ell+1)^3.
\]

Expanding a visible monomial through the hidden marginal produces tuples of
joint configurations.  For each tuple, record how many entries activate each
cubic feature.  Call this its cubic feature profile.  If two collections of
visible monomials have the same multiset of expanded profiles, then every
3-local face--Gibbs representation forces the two corresponding sums of
monomials to be equal.  This remains true when the joint law lies on a
boundary face and has zero cells; this is the boundary-safe marginal-trade
certificate proved in Lean.

There are far more sets of candidates than possible profile histograms once
\[
  n+\ell+3+(n+1)(n+\ell+1)^3 < 2^n-1.       \tag{1}
\]
A finite pigeonhole argument then produces two distinct sets of candidates
with the same profile histogram.  Cubic locality says their monomial sums
must be equal, whereas the nested binary encoding says they cannot be equal.
That contradiction rules out the proposed $\ell$-hidden-bit localization.

The checked arithmetic specialization is
\[
  n=4m+24,
  \qquad
  \ell\le 2^m.
\]
Lean verifies that these values satisfy (1) for every natural number $m$.

\section{The resulting lower bound}

Here $\LC_3(D)$ denotes the least number of hidden binary coordinates in any
cubic localization of $D$.

\begin{theorem}[Lean checked]
For every natural number $m$, the rational distribution $D_m$ above has full
support and satisfies
\[
  \LC_3(D_m)>2^m.
\]
Since $n_m=4m+24$, this can also be written
\[
  \LC_3(D_m)>2^{(n_m-24)/4}=\frac{1}{64}\,2^{n_m/4}.
\]
Moreover, for $m\ge 13$,
\[
  \LC_3(D_m)>n_m^2.
\]
\end{theorem}

Because localization complexity is integer-valued, the first strict
inequality says that at least $2^m+1$ hidden bits are required.  Some sample
parameters are
\[
\begin{array}{rrrr}
\toprule
m & n_m & 2^m & \text{checked conclusion}\\
\midrule
0  & 24 & 1    & \LC_3(D_0)\ge 2\\
1  & 28 & 2    & \LC_3(D_1)\ge 3\\
4  & 40 & 16   & \LC_3(D_4)\ge 17\\
13 & 76 & 8192 & \LC_3(D_{13})\ge 8193\\
\bottomrule
\end{array}
\]
At $m=13$, note that $76^2=5776<8192$.

\section{What ``explicit'' does and does not mean}

The family is explicit in a mathematical sense: given $m$ and $x$, the
formula above determines an exact rational number $D_m(x)$ without referring
to a genericity argument or choosing the distribution by pigeonhole.  It is
also full support, so this result has no direct implication for Boolean
support-circuit lower bounds.

It is not an efficient finite-precision construction.  Near the end of the
table, $\val(x)$ is of order $2^n$, so the binary length of
$2^{2^{\val(x)}}$ is itself doubly exponential in $n$.  The checked
certificates have degree
\[
  2(2^n-1),
\]
and the colliding candidate families are selected noncomputably by the finite
pigeonhole proof rather than printed as practical certificates.

The exponent $1/4$ is also conservative.  The cubic profile count naturally
contains a factor of roughly $(n+\ell)^3$, suggesting that a sharper analysis
of the same mechanism should approach the scale $2^{n/3}$.  What is currently
checked is the clean uniform bound $2^{(n-24)/4}$.

\section{A bounded-precision Boolean tilt and the circuit route}

There is now a second, complementary consequence of the same profile-collision
argument.  Given a Boolean function $f:\bits^n\to\bits$, define
\[
  D_f(x):=\frac{2^{f(x)}}{\sum_y 2^{f(y)}}
  =\frac{1+f(x)}{2^n+|f^{-1}(1)|}.
\]
Unlike $D_m$, this distribution has only two unnormalized weights, $1$ and
$2$.  Its probability table has modest rational precision and full support.

Suppose a sequential NAND circuit with $s$ gates computes $f$.  Lift $D_f$ to
the unique full wire trace of the circuit.  Its support is exactly the graph
of the NAND computation, exposed by the sum of the quadratic NAND penalties.
On that graph, its log density is a constant plus a unary potential on the
output wire.  Therefore the gate wires themselves are $s$ hidden bits of a
quadratic localization, with no additive overhead:
\[
  \LC_2(D_f)\le s,
  \qquad
  \LC_3(D_f)\le s.                         \tag{2}
\]
This circuit-to-localization bridge is now checked in Lean.  Consequently,
any lower bound for $\LC_3(D_f)$ is immediately a lower bound for the NAND
gate complexity of $f$.

The marginal-trade calculation also becomes especially clean.  For a visible
tuple $(x_1,\ldots,x_d)$,
\[
  \prod_{i=1}^d D_f(x_i)
  =Z_f^{-d}\,2^{\sum_i f(x_i)}.
\]
Hence a homogeneous trade detects $D_f$ exactly when the two natural numbers
\[
  \sum_{t\in A}2^{\#\{i:f(x_{t,i})=1\}}
  \quad\hbox{and}\quad
  \sum_{t\in B}2^{\#\{i:f(x_{t,i})=1\}}
                                                        \tag{3}
\]
differ.  The normalizer cancels completely.

The checked powers-of-powers proof chooses two distinct families $A$ and $B$
of binary candidate monomials having the same expanded cubic profile
histogram.  A new interpolation lemma observes that the matrix
\[
  K(T,C)=2^{|T\cap C|}
\]
is invertible over the integers: its one-coordinate matrix is
\[
  \begin{pmatrix}1&1\\1&2\end{pmatrix},
  \qquad
  \begin{pmatrix}1&1\\1&2\end{pmatrix}^{-1}
  =\begin{pmatrix}2&-1\\-1&1\end{pmatrix}.
\]
Tensoring over the candidate coordinates shows that some binary test $T$
separates any distinct $A$ and $B$.  Interpreting $T$ as the truth set of a
Boolean function makes the two quantities in (3) unequal.

This gives the following Lean-checked theorem.

\begin{theorem}[Lean checked, existential]
For every natural number $m$, there exists a Boolean function
\[
  f_m:\bits^{4m+24}\longrightarrow\bits
\]
such that its two-level full-support rational tilt satisfies
\[
  \LC_3(D_{f_m})>2^m.
\]
Every sequential NAND circuit computing $f_m$ therefore has more than $2^m$
gates.
\end{theorem}

This is a real circuit lower bound statement, but it is not an explicit
circuit lower bound.  The function $f_m$ is selected from the noncomputably
chosen profile collision.  In complexity-theoretic terms, this is an
existence/counting lower bound in a new language, not a breakthrough against
explicit circuit lower bounds.

The remaining target is now unusually crisp: replace the interpolated test
$T$ by a uniformly and efficiently describable test---for example a
code-theoretic, algebraic, or pseudorandomly generated subset---while
retaining separation of the matched profile families.  Accomplishing that
with $m=\Omega(\log n)$ would already give a superlinear explicit NAND lower
bound through (2), so one should expect this step to contain the genuine
circuit-lower-bound difficulty.

\section{The exact linear-algebra bottleneck}

There is a useful way to state that remaining difficulty without hiding it
behind the pigeonhole argument.  Fix a hidden-bit budget $\ell$.  Index the
binary candidates by $C$, and let $A_\ell(C)$ be the integer histogram of all
expanded cubic feature profiles contributed by candidate $C$.  A signed
integer vector $b=(b_C)_C$ is a homogeneous marginal trade precisely when
\[
  A_\ell b=0.
\]
Its evaluation on the Boolean test $T$ is
\[
  K_Tb=\sum_C b_C2^{|T\cap C|}.
\]
Thus the concrete certificate problem for an explicit function is
\[
  A_\ell b=0,\qquad K_Tb\ne0.                 \tag{4}
\]
Over the rationals, such a $b$ exists exactly when the row $K_T$ does not lie
in the row span of $A_\ell$.  Clearing denominators turns a rational witness
into an integer trade.

The tensor inverse proves that the full matrix $K$ has full rank.  Therefore
some row $K_T$ escapes every proper row span, which is exactly the existential
theorem above.  Full rank alone says nothing about the description or circuit
complexity of the escaping row.  This is why interpolation closes the
counting theorem but not the explicit theorem.

There is also a sharp sanity check.  If a proposed test $T$ is computed by an
$s$-gate NAND circuit, the checked trace construction gives an $s$-hidden-bit
localization of $D_T$.  Consequently (4) cannot hold at any budget
$\ell\ge s$.  In particular, parity, affine tests, and other tests with known
linear-size circuits cannot possibly yield a superlinear bound by a clever
choice of trade.  A successful explicit $T$ must be a function for which the
row-span nonmembership proof is itself doing genuine circuit-lower-bound
work.

This suggests a focused research program rather than a blind search: expose
enough symmetry in $A_\ell$ to diagonalize part of its row space (for example
with code or association-scheme methods), and then place an explicit $K_T$
in a missing component.  A pseudorandom hitting set would help only after one
proves additional structure---such as bounded degree, sparsity, or Fourier
concentration---for the relevant trade kernel; arbitrary nonzero multilinear
polynomials can vanish on all but one Boolean test.

\section{A structured Fourier laboratory}

The most useful structured family found so far replaces the arbitrary
candidate fiber by a doubly exponential one with an exact Fourier transform.
Let
\[
  A=\bits^q,\qquad Z=\bits^4,\qquad N=|A|=2^q.
\]
Use one additional marker coordinate, fixed to zero.  For each truth table
$s:A\to\bits$, put
\[
  C_s:=\{(0,a,z):\operatorname{parity}(z)=s(a)\}.
\]
Every $C_s$ has $8N$ points.  More importantly, all $2^N$ sets $C_s$ have
the same visible moments through degree three.  Indeed, after fixing any at
most three of the four $z$ coordinates, one $z$ coordinate remains free and
is uniquely determined by the requested parity.

If the test is another member $C_t$ of the family, then
\[
  |C_s\cap C_t|=8\,|\{a:s(a)=t(a)\}|,
\]
and hence the Boolean-tilt response matrix is
\[
  K(t,s)=256^{N-d_H(s,t)}.
\]
It is the $N$-fold tensor power of
\[
  \begin{pmatrix}256&1\\1&256\end{pmatrix}.
\]
The Walsh character indexed by $U\subseteq A$ has eigenvalue
\[
  257^{N-|U|}255^{|U|},
\]
so every eigenvalue is nonzero.  This calculation is now checked in Lean.
Lean also checks, uniformly in $q$ and in the two truth tables, that all of
the sets $C_s$ have the same cubic moment profile.
It is a clean diagonalization of the test side of the problem, but it also
rules out a tempting shortcut: the test kernel itself has no missing Fourier
level.

For one hidden bit, after removing a common visible monomial, let $b_{a,z}$
denote the relevant quadratic toric monomial and set
\[
 E_{a,e}:=\prod_{z:\operatorname{parity}(z)=e}(1+b_{a,z}).
\]
The column for $s$ is
\[
  P_s=\prod_{a\in A}E_{a,s(a)}.
\]
Its Walsh transform factors exactly:
\[
 \sum_s(-1)^{\sum_{a\in U}s(a)}P_s
 =\prod_{a\notin U}(E_{a,0}+E_{a,1})
  \prod_{a\in U}(E_{a,0}-E_{a,1}).              \tag{5}
\]
Lean now checks this identity over an arbitrary commutative coefficient
ring.  Thus the remaining efficient-explicitness problem is not about an
opaque matrix: it is about the span of the explicit toric products on the
right side of (5).

There is also a useful triangular filtration.  Giving each $b_{a,z}$ degree
one,
\[
 E_{a,0}+E_{a,1}=2+O(b),\qquad
 E_{a,0}-E_{a,1}=g_a+O(b^2),
\]
where
\[
 g_a=\sum_z(-1)^{\operatorname{parity}(z)}b_{a,z}.
\]
Therefore the least homogeneous degree of the $U$th Walsh product is $|U|$,
with leading part
\[
  2^{N-|U|}\prod_{a\in U}g_a.                  \tag{6}
\]
Any relation among (5) must begin with a relation among the simpler products
in (6).  For a quadratic toric model one can further write
\[
 b_{a,z}=c_a d_z\prod_{r=1}^4u_r(a)^{z_r},
\]
so, after the quadratic prefix factor $c_a$ is removed, $(g_a)_a$ is a sum
of at most sixteen product tensors.  Flattening minors are consequently a
real lead.  The difficulty is equally precise: $c_a$ contains arbitrary
pair interactions among prefix bits, and its entrywise product can restore
full flattening rank.  A successful invariant has to cancel the whole
quadratic prefix profile.

\section{Two honest no-go results}

First, the much larger symmetric group on the $N$ blocks is not a symmetry
of the binary-prefix construction.  Its literal block symmetries are only
those induced by signed permutations of the $q$ prefix coordinates.  Treating
the model as the full Hamming association scheme would therefore prove a
statement about the wrong matrix.

One can make $S_N$ literal by encoding block $a$ with a one-hot prefix.
Unfortunately that gives the cubic joint model private parameters for every
block.  Already for one hidden bit it admits the specialization
\[
 b_{a,z}=u_a\prod_{r=1}^4v_{a,r}^{z_r}.
\]
The two polynomials $E_{a,0}$ and $E_{a,1}$ are linearly independent: their
constant terms agree, while the coefficient of $u_a$ in their difference is,
up to sign,
\[
  \prod_{r=1}^4(1-v_{a,r})\ne0.
\]
Variables belonging to different blocks are disjoint, so tensoring these
independent pairs proves that all columns $P_s$ are independent.  Thus the
one-hot version has the desired symmetry and no trade at all.  There is a
genuine tension between large block symmetry and the private local parameters
that symmetry introduces.

Second, a familiar easy test cannot be rescued by a more ingenious trade.
The Lean-checked circuit trace theorem says that if $T$ has a NAND circuit
with at most $\ell$ gates, then every $\ell$-hidden marginal trade evaluates
to zero on $D_T$.  Parity, affine functions, and other linear-size tests are
therefore categorically incapable of proving a superlinear bound in this
framework.

The exact target can be written without euphemism.  For a truth table
$t_q:A\to\bits$ and a budget $L(q)$ growing superlinearly in the visible
dimension $q+5$, construct an integer vector $b_{q,L}$ such that
\[
  M_{q,L}b_{q,L}=0,
  \qquad
  \sum_s b_{q,L}(s)\,256^{N-d_H(s,t_q)}\ne0.    \tag{7}
\]
The first equation is a complete cubic-profile cancellation, including all
hidden assignments and all boundary faces.  The second says that the chosen
Boolean tilt detects it.

The literal finite problem is now solved.  Set $L(q)=q^2$.  For every
$q\ge64$, compare all subsets of the $2^N$ candidate columns with their
complete expanded cubic-profile histograms.  The number of subsets is larger
than the number of histograms, so two distinct subsets collide.  Choose the
lexicographically first ordered collision $(A_q,B_q)$ and put
\[
  b_q(s)={\bf1}_{A_q}(s)-{\bf1}_{B_q}(s).
\]
Then $b_q\ne0$ and $M_{q,q^2}b_q=0$.  Since the $256$-agreement tensor is
integer-injective, some row detects $b_q$; choose the numerically first such
row as $t_q$.  This gives (7) with $L(q)=q^2=\omega(q+5)$.

The same collision compiles to a boundary-safe face--Gibbs marginal-trade
certificate.  Consequently the two-level rational law
\[
  D_q(x)\ \propto\ 2^{{\bf1}_{C_{t_q}}(x)}
\]
on $q+5$ visible variables has full support and
$\LC_3(D_q)>q^2$.  Every cell has one of only two unnormalized weights, so
the lower bound is genuinely bounded-precision.

This does \emph{not} constitute an efficiently explicit circuit lower bound.
Both $(A_q,B_q)$ and $t_q$ are defined by finite exhaustive searches over
doubly exponential objects.  In ordinary complexity terminology this is a
canonical diagonalization family, not a low-description or efficiently
computable family.  The ambitious remaining target is to replace those
first-witness searches by a uniform efficient construction while retaining
(7).  That strengthened statement would contain genuine general-circuit
lower-bound content.

\section{Lean correspondence}

The powers-of-powers construction is in
\texttt{KLocality/UniformExplicitCubicLowerBound.lean}; the structured
bounded-precision construction is in the modules beginning with
\texttt{KLocality/BlockParity}.  The main names are
\begin{itemize}
  \item \texttt{uniformCubicCellExponent} and
  \texttt{uniformCubicUnnormalizedRat} for $W_m$;
  \item \texttt{uniformCubicDistribution} for the normalized rational table;
  \item \texttt{superlinearCubicVisibleBits} and
  \texttt{superlinearCubicDistribution} for $n_m$ and $D_m$;
  \item \texttt{superlinearCubicDistribution\_support} for full support;
  \item \texttt{superlinearCubicDistribution\_localizationComplexity\_gt}
  for $\LC_3(D_m)>2^m$;
  \item
  \texttt{superlinearCubicDistribution\_localizationComplexity\_gt\_sq}
  for the quadratic corollary.
  \item \texttt{booleanTiltDistribution} for $D_f$;
  \item
  \texttt{localizationComplexityBits\_three\_booleanTilt\_le\_CNAND}
  for the circuit-to-localization bridge;
  \item \texttt{binarySubsetFamilyProfile\_injective} for the interpolation
  lemma;
  \item
  \texttt{exists\_superlinearBooleanTilt\_}\\
  \texttt{nandGateLowerBound}
  for the existential NAND lower bound.
  \item
  \texttt{booleanTiltCodes\_eq\_of\_nandCircuit\_le}
  for the simple-circuit sanity check.
  \item \texttt{binaryAgreementKernel\_walsh\_eigenvector} for the exact
  Walsh spectrum of the block-parity test kernel;
  \item \texttt{binaryProductColumn\_walshTransform} for (5);
  \item \texttt{parityFourClass\_cubicMoment\_eq} and
  \texttt{two\_pow\_parityFourClass\_inter\_card} for the local four-bit
  parity-block facts;
  \item \texttt{blockParityMoment\_eq} for the full arbitrary-width common
  cubic-moment fiber.
  \item \texttt{blockParityCanonicalTradeWitness} for the exact nonzero
  $b_q$, matrix-kernel identity, and detecting $t_q$ at $L(q)=q^2$.
  \item \texttt{blockParityCanonicalCertificate} for the compiled
  boundary-safe marginal trade.
  \item
  \texttt{blockParityCanonicalDistribution\_}\
  \texttt{localizationComplexity\_gt} for
  $\LC_3(D_q)>q^2$ on a visible type of cardinality $q+5$.
\end{itemize}

\end{document}
